\documentclass[10pt,a4paper]{article}
\usepackage[utf8]{inputenc}
\usepackage[a4paper,margin=0.8in]{geometry}
\usepackage{graphicx}
\usepackage{amsmath,amssymb}
\usepackage{times}
\usepackage{multicol}
\usepackage{ragged2e}
\usepackage{caption}
\usepackage{subcaption}
\usepackage{array}
\usepackage{booktabs}
\usepackage{float}
\usepackage{titlesec}
\usepackage{setspace}
\usepackage{hyperref}
\usepackage{xcolor}
\hypersetup{colorlinks=true, citecolor=blue}
\setlength{\columnsep}{0.8cm}
\titleformat{\section}
{\normalfont\large\bfseries}
{\thesection.}{0.5em}{}

\titleformat{\subsection}
{\normalfont\large\bfseries}
{\thesubsection}{0.5em}{}

\titleformat{\subsubsection}
{\normalfont\normalsize\bfseries}
{\thesubsubsection}{0.5em}{}
\captionsetup{
font=small,
labelfont=bf,
justification=centering
}
\numberwithin{figure}{section}
\numberwithin{table}{section}
\begin{document}
\twocolumn[
\begin{@twocolumnfalse}
\noindent
{\Large \textbf{Development and Characterization of Heat Treatable Hybrid Metal Matrix Composites}}

\vspace{0.2cm}

\noindent
{\normalsize  \textbf{ K.AKASH(1GC24IS009)}}

\vspace{0.2cm}

\noindent
{\footnotesize
Department of Information Science and Engineering, Ghousia College of Engineering, BM Road Bengaluru south Ramanagara.,\\
Visvesvaraya Technological University, Karnataka, India-562159\\
akxshh28@gmail.com
}

\vspace{0.3cm}

\hline

\vspace{0.3cm}

\noindent
{\normalsize \textbf{ABSTRACT}}

\vspace{0.1cm}

\justify
\footnotesize
 In recent years, there has been an ever-increasing demand for enhancing mechanical properties of Aluminum Matrix Composites (AMCs), which are finding wide applications in the field of aerospace, automobile, defence etc. Among all available aluminium alloys, Al6061 is extensively used owing to its excellent wear resistance and ease of processing. Newer techniques of improving the hardness and wear resistance of Al6061 by dispersing an appropriate mixture of hard ceramic powder and whiskers in the aluminium alloy are gaining popularity. The conventional aluminium based composites possess only one type of reinforcements. Addition of hard reinforcements such as silicon carbide, alumina, titanium carbide, improves hardness, strength and wear resistance of the composites. However, these composites possessing hard reinforcement do possess several problems during their machining operation. AMCs reinforced with particles of Graphite have been reported to possess better wear characteristics owing to the reduced wear because of formation of a thin layer of Gr particles, which prevents metal to metal contact of the sliding surfaces. Further, heat treatment has a profound influence on mechanical properties of heat treatable aluminium alloys and its composites. For a solutionising temperature of 530$^\circ$C, solutionising duration of 1hr, ageing temperature of 175$^\circ$C, quenching media and ageing duration significantly alters mechanical properties of both aluminium alloy and its composites.In the light of the above, the present paper aims at developing aluminium based hybrid metal matrix composites containing both silicon carbide and graphite and characterize their mechanical properties by subjecting it to heat treatment.

\vspace{0.3cm}
\footnotesize
\noindent
\textit{Keywords: Hybrid Metal Matrix Composites, Microhardness, Heat treatment, Vortex casting technique.}

\vspace{0.5cm}

\hline

\vspace{0.6cm}

\end{@twocolumnfalse}
]
\section{Introduction}

\justify
\large
In recent years aluminium matrix composites) are gaining widespread popularity in several technological sectors owing to their excellent  corrosion and wear resistance, higher fatigue life, good high temperature oxidation resistance in addition to being light in weight when compared with conventional alloys. At present AMCs are attractive alternatives for aerospace and automotiv  applications because of their high stiffness-to-weight characteristics   \cite{rar1}  . Currently, focus on development of aluminium, copper, magnesium, titanium based metal matrix composites is carried out to explore their possible applications in several high-tech areas. The various reinforcements that have been tried out to develop AMCs are graphite, silicon carbide, titanium carbide, tungsten, boron, Al$_2$O$_3$, flyash, Zr, Si$_3$N$_4$, TiB$_2$. The conventional aluminium based composites possess only one type of reinforcements. Addition of  hard reinforcements such as silicon carbide, alumina, titanium carbide, improves hardness, strength and wear resistance of the composites \cite{rar4}.  These composites possessing hard reinforcement do posses several problems during their machining operation. AMCs reinforced with particles of Gr have been reported to be possessing better wear characteristics owing to the reduced wear because of formation of a thin layer of Gr particles, which prevents metal to metal contact of the sliding surfaces. AMCs reinforced with SiC particulates are known for higher modulus, strength and wear resistance compared to conventional alloys. Addition of SiC particulates increases both mechanical strength and wear resistance of Al alloy. But the consequent increase in hardness makes the machining difficult. On the other hand, addition of Gr particulates facilitates easy machining and results in reduced wear of Al–Gr composites compared to Al alloy \cite{rar1}. It is reported that the surface finish of the hard reinforced metal matrix composites are inferior when compared with the matrix alloy. Further it is absorbed that during turning, the hard reinforced metal matrix composites resulted in higher flank wear with increased content of the reinforcement. It is reported that composites possessing softer reinfor- cement - possess good machinability index \cite{rar10}.  Hence the current interest is to produce Hybrid Metal Matrix Composites were in more than one type, shape and size of the reinforcement are used to obtain synergistic properties of the reinforcement and the matrix chosen. It is reported that hybridization of  reinforcement has enhanced structural, physical, mechanical and tribological behavior of HMMC's  when compared with metal matrix alloy\cite{rar5}. Hence attempts are made to develop aluminum based hybrid metal matrix composites consisting of both hard reinforcement (SiC) and soft reinforcement (Gr) in the present paper. Further to enhance the mechanical properties, the specimens are subjected to heat treatment processes.


\section{Experimental details}



\subsection{Composite preparation}

Al6061 based composites were prepared by vortex method of liquid metallurgy route. A quantity of 3kgs of Al6061 alloy was used each time in an electric melting furnace with graphite crucible for melting with furnace temperature set at 710$^\circ$C. Silicon carbide particles of 10 micron size and graphite particles of 60 micron size were used. The permanent molds of cast iron along with the reinforcements were heated in order to reduce the effect of chilling during solidification \cite{rar7,rar6}   . Degassing of the melt was done with commercially available tablets of hexachloroethane (C$_2$Cl$_6$). After degassing, the preheated SiC and Gr were added slowly into the vortex while continuing the stirring process up to 10 minutes. The amount of reinforcement was varied from 1wt\% to 4wt\% in steps of 1wt\% of Gr keeping constant 7wt\% SiC.



\subsection{Evaluation of hardness}

Vickers micro hardness test was performed on all samples of Al6061 and its hybrid composites for both with and without heat treatment cases. The polished samples were subjected for micro hardness tests using shimadzu micro hardness tester as shown in Fig. 2.2. A load of 100 g for period of 10 sec was applied. The hardness was noted by taking the diagonal lengths of indentation produced.
The tests were carried out at five different locations in order to negate the possible effect of indenter resting on the harder particles. The average of all the five readings was taken as hardness of sample.



\begin{figure}[H]
\centering
\includegraphics[width=0.50\columnwidth]{tes}
\caption{Photograph of Shimadzu microhardness tester}
\end{figure}



\begin{table}[H]
\centering
\caption{Specifications of shimadzu microhardness tester}
\renewcommand{\arraystretch}{1.1}
\begin{tabular}{|p{4cm}|p{5cm}|}
\hline
\textbf{Parameter} & \textbf{Specification}
 \\ \hline
Make & Shimadzu, Japan 
\\ \hline
Loading method & Lever method by Electric automatic loading system
 \\ \hline
Period of loading & 5, 10, 15, 30, \& 45 sec  
\\ \hline
Loading weight & 15, 25, 30, 100, 150, 200, 300, 500, 800 \& 1000 gram
 \\ \hline
Indenter & Vickers
 \\ \hline
Magnifications & 100 times
\\ \hline 
Microscope & 400 times 
\\ \hline
Measuring objectives & 40 times  
\\ \hline
Max. measuring scale & 200$\mu$m
 \\ \hline
Standard scale & With one division of 20$\mu$m
 \\ \hline
Measuring scale & With one division of 0.5$\mu$m 
\\ \hline
\end{tabular}

\end{table}


\section{Results and discussion}



\subsection{Scanning electron micrograph studies}

The scanning electron micrographs of Al6061 and its composites are shown in Fig. 3.1(a-f).


\begin{figure}[H]
\centering
\begin{subfigure}{0.34\textwidth}
\centering
\includegraphics[width=1.25\linewidth]{al601}
\centering 
(a)
\end{subfigure}
\hfill
\begin{subfigure}{0.34\textwidth}
\centering
\includegraphics[width=1.25\linewidth]{al602}
\centering 
(b)
\end{subfigure}
\end{figure}

\centering
\begin{figure}[H]
\centering
\begin{subfigure}{0.35\textwidth}
\centering
\includegraphics[width=1\linewidth]{al6011}
\centering 
(c)
\end{subfigure}
\hfill
\begin{subfigure}{0.33\textwidth}
\centering
\includegraphics[width=1\linewidth]{al6022}
\centering
(d)
\end{subfigure}
\end{figure}

\begin{figure}[H]4
\centering
\begin{subfigure}{0.35\textwidth}
\centering
\includegraphics[width=1\linewidth]{al603}
\centering 
(e)
\end{subfigure}
\hfill
\begin{subfigure}{0.36\textwidth}
\centering
\includegraphics[width=1\linewidth]{al604}
\centering
(f)
\end{subfigure}

\caption{Optical microphotographs of Al6061 and its hybrid composites}

\end{figure}
\justify
It is observed from Fig. 3.1(a-f) that the SiC and Gr particles are fairly uniformly distributed. The extent of porosity noticed is also less. There is a clear evidence of homogeneous distribution of reinforcements, although a clear cut identification of SiC and Gr is quite difficult.


\subsection{Hardness test}

\subsubsection{Effect of reinforcements without heat treatment}

The effect of reinforcements on microhardness of Al6061 alloy and its composites is as shown in Fig. 3.2.



\begin{figure}[H]
\centering
\includegraphics[width=0.9\columnwidth]{wih}
\caption{Effect of reinforcements on microhardness of Al6061 alloy and its composites}
\end{figure}
\justify
It is observed from Fig. 3.2 that the hardness of the Al6061 increases with addition of silicon carbide. But with addition of graphite the hardness decreases.
To overcome this, the SiC hard ceramic particle wasadded and was maintained constant 7wt\%. Sicparticles can act as the obstacles to the movement of dislocation as referred by Mahalingegowda and T Rajmohan  \cite{rar8,rar9}.   The SiC particles in the matrix alloy provide protection to the softer matrix. Thus, limiting 
the deformation and also resists the penetration and cutting of slides on the surface of the composites. Hardness of all the hybrid composites was 
significantly greater than that of the base alloy characterized to the hard nature of SiC particles.The higher hardness values for the hybrid composites 
containing 7wt\% of SiC is due to the presence of hard SiC particles. This result is a good agreement with the result of F. Akhlaghi and Zare-bidaki   \cite{rar2}. 


\subsubsection{Effect of reinforcements with heat treatment}

The effect of reinforcements with variation of ageing time on microhardness of heat treated Al6061 alloy and its composites is shown in Fig. 3.3(a-c).



\begin{figure}[H]
\centering

\begin{subfigure}{1\columnwidth}
\centering
\includegraphics[width=\linewidth]{air}
\caption{}
\end{subfigure}

\vspace{0.5cm}

\begin{subfigure}{1.0\columnwidth}
\centering
\includegraphics[width=\linewidth]{wat}
\caption{}
\end{subfigure}

\vspace{0.5cm}

\begin{subfigure}{1.0\columnwidth}
\centering
\includegraphics[width=\linewidth,height=5cm,keepaspectratio]{ice}
\caption{}
\end{subfigure}

\caption{Effect of reinforcements with variation of ageing time on microhardness of heat treated Al6061 alloy and its composites}

\end{figure}

It is observed from Fig. 3.3(a-c) that the hardness increases with heat treatment.
Heat treatment has a profound influence on the microhardness of the matrix alloy as well as its hybrid composites. For a solutionising temperature of 530$^\circ$C, for duration of 1hr, ageing temperature of 175$^\circ$C, quenching media and ageing duration significantly enhances the microhardness of both the matrix alloy and its composites.


\section{Conclusion}
Al6061-SiC-Gr hybrid composites have been successfully produced by vortex method upto 4wt\% Gr with constant 7wt\% SiC. Fabrication of hybrid composites with more than 4wt\% of Gr with constant 7wt\% SiC was not achieved successfully due to low density of graphite. Scanning electron micrograph studies clearly revealed uniformity in the distribution of reinforcements and excellent bond between the matrix and the reinforcement. Microhardness of Al6061 decreases with increase in graphite content. Microhardness of Al6061 based hybrid composites is higher when compared with that of the matrix alloy . Presence of hard silicon carbide reinforcement in the hybrid composites leads to enhancement in microhardness of hybrid composites. Further heat treatment has a profound influence on the microhardness of the matrix alloy as well as its hybrid composites. For all the heat treatment processes studied, ice quenching with ageing duration of 6hrs resulted in improved hardness of both the unreinforced matrix alloy and its hybrid composites.
\clearpage
\renewcommand{\bibname}{Reference}
\bibliographystyle{unsrt}
\bibliography{rar}
\end{document}